\documentclass[11pt,a4paper]{article}

\usepackage[utf8]{inputenc}
\usepackage[T1]{fontenc}
\usepackage{geometry}
\usepackage{hyperref}
\usepackage{enumitem}
\usepackage{xcolor}
\usepackage{fancyhdr}
\usepackage{amsmath}
\usepackage{amssymb}
\usepackage{booktabs}
\usepackage{listings}

\geometry{margin=1in}
\hypersetup{
  colorlinks=true,
  linkcolor=blue,
  urlcolor=blue
}

\setlist{itemsep=0.35em, topsep=0.35em}

\pagestyle{fancy}
\fancyhf{}
\lhead{Elements of AIML (B.Tech)}
\rhead{Unit 02: Diagnostic Assessment}
\cfoot{\thepage}
\setlength{\headheight}{14pt}

\lstset{
  basicstyle=\ttfamily\small,
  keywordstyle=\color{blue},
  commentstyle=\color{gray},
  stringstyle=\color{teal},
  showstringspaces=false,
  columns=fullflexible,
  frame=single
}

\newcommand{\SubmissionDeadline}{March 8, 2026}

% Marks per question (total = 100)
\newcommand{\EMarks}{\textbf{[2 marks]} }
\newcommand{\MMarks}{\textbf{[3 marks]} }
\newcommand{\HMarks}{\textbf{[5 marks]} }

\begin{document}

\begin{center}
  {\LARGE \textbf{Diagnostic Assessment -- Unit 02}}\\[0.25em]
  {\large Elements of AIML}\\[0.25em]
  \normalsize (Logic for AI; designed to identify learning gaps)
\end{center}

\noindent
\textbf{Instructor:} Tofik Ali \hfill \textbf{Time Limit:} None\\
\textbf{Submission Deadline:} \SubmissionDeadline\\
\textbf{Student Name:} \_\_\_\_\_\_\_\_\_\_\_\_\_\_\_\_\_ \hfill \textbf{Roll No.:} \_\_\_\_\_\_\_\_\\

\vspace{0.6em}
\hrule
\vspace{0.8em}

\textbf{Instructions}
\begin{itemize}[leftmargin=*]
  \item \textbf{This is a diagnostic assessment, not a quiz/test.} The goal is to identify learning gaps so you know what to revise next.
  \item \textbf{Total marks: 100.} Easy: 10 $\times$ 2 = 20, Medium: 10 $\times$ 3 = 30, Hard: 10 $\times$ 5 = 50.
  \item \textbf{No time limit.} Submit on or before \textbf{\SubmissionDeadline}.
  \item \textbf{Important (70\% rule): To have your assignment marks counted, you must score at least 70\% (70/100 or more) in this assessment.}
    If you score below 70\%, you may reattempt and resubmit. Your latest submission will be considered.
  \item \textbf{Marking policy (honest attempt):} Marks are deducted only for (i) not attempting a question and (ii) high similarity with other submissions (copying).
    Correctness is \textbf{not} the main focus; your reasoning and effort matter more.
  \item \textbf{Show your real effort:} if you tried multiple approaches for any question, write them all (Attempt 1, Attempt 2, ...). Partial/incorrect attempts are acceptable.
  \item \textbf{Many questions are open-ended by design:} multiple solutions/variants are acceptable. Use your own example data and explain your choices.
  \item Unless specified, you may choose your own symbols/assumptions (e.g., domain, predicate names), but you must clearly mention your choices.
  \item Submission format: submit a single ZIP containing \texttt{notebook.ipynb} and \texttt{README.md} (plus optional \texttt{assets/}).
  \item Show key steps for logic problems (truth tables, CNF steps, resolvents, substitutions).
  \item Round final numeric answers to \textbf{2 decimal places} (when applicable).
\end{itemize}

\vspace{0.6em}

\section*{Easy (10 Questions)}
\begin{enumerate}[label=\textbf{E\arabic*.}, leftmargin=*]
  \item \EMarks \textbf{[Subjective]} Define a \textbf{model} in propositional logic (2--4 lines).

  \item \EMarks \textbf{[MCQ -- Select one]} A sentence is \textbf{valid} if it is:
    \begin{enumerate}[label=(\Alph*), leftmargin=*]
      \item true in at least one model
      \item true in all models
      \item false in all models
      \item neither true nor false
    \end{enumerate}

  \item \EMarks \textbf{[Subjective]} Write the truth table for $P \rightarrow Q$ and state when it is false.

  \item \EMarks \textbf{[Subjective]} Explain (in 4--6 lines): $KB \models \alpha$ vs $KB \rightarrow \alpha$.

  \item \EMarks \textbf{[Subjective]} Convert $P \rightarrow Q$ into an equivalent expression using only $\lnot$ and $\lor$.

  \item \EMarks \textbf{[Subjective]} Convert $\lnot(P \rightarrow Q)$ into CNF.

  \item \EMarks \textbf{[Subjective]} Perform one resolution step on:
    \[
      (P \lor Q),\ (\lnot P \lor R).
    \]

  \item \EMarks \textbf{[MCQ -- Select all that apply]} Which symbols are \textbf{quantifiers} in FOL?
    \begin{enumerate}[label=(\Alph*), leftmargin=*]
      \item $\forall$
      \item $\exists$
      \item $\land$
      \item $\rightarrow$
      \item $\lor$
    \end{enumerate}

  \item \EMarks \textbf{[Subjective]} In $Teaches(Prof, Course)$ identify predicate symbol, constants, and variables (if any).

  \item \EMarks \textbf{[Subjective]} State the difference between a \textbf{function symbol} and a \textbf{predicate symbol} (4--6 lines).
\end{enumerate}

\section*{Medium (10 Questions)}
\begin{enumerate}[label=\textbf{M\arabic*.}, leftmargin=*]
  \item \MMarks \textbf{[Subjective]} Let $KB=\{P \rightarrow Q,\ P\}$.
    Check whether $KB \models Q$ using truth tables or model checking.

  \item \MMarks \textbf{[Subjective]} Determine whether the set is satisfiable:
    \[
      (P \lor Q),\ (\lnot P \lor Q),\ (P \lor \lnot Q).
    \]
    Provide one satisfying assignment if it exists.

  \item \MMarks \textbf{[Subjective]} Convert to CNF:
    \[
      (P \rightarrow Q) \land (Q \rightarrow R).
    \]

  \item \MMarks \textbf{[Subjective]} Use resolution refutation to show:
    \[
      \{P \rightarrow Q,\ Q \rightarrow R,\ P\} \models R.
    \]
    (Show clauses and at least the main resolution steps.)

  \item \MMarks \textbf{[Subjective]} Translate into FOL (domain: students and courses):
    ``Every student studies at least one course.''

  \item \MMarks \textbf{[Subjective]} Explain (6--8 lines) why we convert formulas into CNF before applying resolution.

  \item \MMarks \textbf{[Subjective]} Find an MGU (if it exists):
    \[
      Knows(John, x)\ \text{and}\ Knows(y, Bill).
    \]

  \item \MMarks \textbf{[Subjective]} Explain occurs-check with the example $x$ vs $f(x)$ (6--8 lines).

  \item \MMarks \textbf{[Subjective]} Resolve with unification:
    \[
      (P(x) \lor Q(x)),\ (\lnot P(a)).
    \]
    Show the substitution used.

  \item \MMarks \textbf{[Subjective]} Explain in 8--10 lines: why propositional logic is not enough for representing relations like ``Parent(x,y)''.
\end{enumerate}

\section*{Hard (10 Questions)}
\begin{enumerate}[label=\textbf{H\arabic*.}, leftmargin=*]
  \item \HMarks \textbf{[Subjective]} Prove using resolution refutation:
    \[
      \{A \rightarrow B,\ B \rightarrow C,\ C \rightarrow D,\ A\} \models D.
    \]

  \item \HMarks \textbf{[Subjective]} Convert to CNF:
    \[
      \lnot(P \leftrightarrow Q) \rightarrow (R \lor P).
    \]
    Show key steps.

  \item \HMarks \textbf{[Subjective]} Convert to clause form (show Skolemization):
    \[
      \forall x\ \exists y\ \forall z\ (Loves(x,y) \rightarrow Loves(z,y)).
    \]

  \item \HMarks \textbf{[Subjective]} Using resolution, prove $Mortal(Socrates)$ from:
    \[
      \forall x\ (Human(x) \rightarrow Mortal(x)),\quad Human(Socrates).
    \]

  \item \HMarks \textbf{[Subjective]} Write two FOL translations for:
    ``Every student likes a course.''
    Explain the difference in meaning (10--12 lines).

  \item \HMarks \textbf{[Subjective]} Perform a unification trace:
    \[
      P(g(x), f(y), y)\ \text{and}\ P(g(a), f(z), z).
    \]

  \item \HMarks \textbf{[Subjective]} Construct an unsatisfiable clause set (at least 5 clauses) and derive the empty clause by resolution.

  \item \HMarks \textbf{[Subjective]} Give a small KB and show a countermodel to demonstrate $KB \not\models \alpha$.
    (Write the assignment/interpretation clearly.)

  \item \HMarks \textbf{[Subjective]} Build a small propositional KB for a security system (alarm/burglary/earthquake)
    with at least 4 rules/facts and ask one entailment query. Solve using models or resolution.

  \item \HMarks \textbf{[Subjective]} Short reflective question (logic in AI):
    In 12--16 lines, explain how entailment, satisfiability, and resolution connect to automated reasoning in AI systems.
\end{enumerate}

\vfill
\noindent\textit{End of Assessment}

\end{document}
